\documentclass[10pt,a4paper]{article}
\usepackage[margin=20mm]{geometry}
\usepackage[T1]{fontenc}
\usepackage{hyperref}
\setlength{\parindent}{0pt}
\setlength{\parskip}{6pt}
\title{OpenCV and Apertus pipeline migration}
\author{See It, Say It, Pick It}
\date{16 October 2026}
\begin{document}
\maketitle
\textbf{Status.} User approved execution. CPU regression checks and the full development-set
mock run passed. Both stacking variants completed without invalid answers. The measured
comparison below uses the official development scorer; the jury score is excluded.
\begin{center}
\small
\begin{tabular}{lrrrrr}
Method & Total /65 & Grounding /35 & Safety /30 & Targets /40 & Seconds \\
\hline
Previous pipeline & 41.6 & 25.4 & 16.2 & 29 & 104 \\
Combined, geometry & 46.0 & 29.8 & 16.2 & 34 & 163 \\
Combined, vision & 45.5 & 28.9 & 16.7 & 33 & 176 \\
\end{tabular}
\end{center}

These are single complete runs via the CSCS Apertus endpoint, with perception rerun each time.
The small difference between stacking methods is within observed run-to-run variation.
Geometry remains the default; vision stacking is available as a configuration override.

\textbf{Source.} Adapted reusable methods and prompts from
\texttt{origin/dev/joswin-opencv}, revision
\texttt{2c1d170}.
The existing LangGraph architecture, calibration, object-region detection and numbered image
annotations remain the integration points. The sole active configuration is
\texttt{configs/pipeline.yaml}; historical manifests keep their original embedded settings.

\textbf{Perception and scene graph.} Apertus names each cyan-outlined region and supplies up to
two plausible alternative names. Names without a colour receive the measured colour. Merging
same-kind regions preserves aliases. The object list includes tray membership, calibrated
size, graspability, image order, stacking, and a complete nearest-to-farthest ranking that
excludes its anchor. The default stack method measures mask overlap against convex hulls;
our Apertus close-up method is retained behind \texttt{stack\_method: vision}.

\textbf{Reasoning and execution checks.} A separate Apertus safety node describes requested
handling before planning. Its failure produces an explicit clarification request after the
configured retry count. Apertus then proposes object/destination moves constrained to detected
ids and known places. Code rejects unknown ids, self-placement, nonfinal holds, too-wide picks,
and handling a detected human hand, and returns the problem to Apertus. Blocking objects are
cleared recursively; the resulting first pick determines the output pixel. Stack cycles are
rejected. New placements update the stack simulation. Requested moves are never silently
removed for graspability. Complete replies missing closing brackets are recovered only when
the entire resulting JSON satisfies its schema.

\textbf{Protocol and reproducibility.} Every run stores the effective YAML configuration,
scene descriptions, raw model replies, retry reasons, results, and a manifest. Generation
settings, prompts, thresholds and optional methods are declared in the single configuration.
Local vLLM remains the default endpoint. A hosted development benchmark is not evidence of a
fully local submission; the final submission must be rerun with a local model endpoint.

\textbf{Scope and limitations.} Ground-truth scene graphs and pre-annotated challenge images
are excluded from runtime. Joswin's scores above 56 in his progress log used ground-truth
perception and are diagnostic ceilings, not photo-only results. Instruction-wording regex hints
are excluded: Apertus alone interprets the instruction. Region thresholds are inherited,
including development-tuned warmth, fragment-area and attached-thickness cutoffs; their
held-out generalisation remains unverified. Table-plane size estimates are biased for tall
objects. A model-assigned plate or saucer label also marks an object ungraspable, following the
challenge rules; a mistaken label can therefore reject a valid pick. Nearest rankings use
interior-point distances on the table plane. Hull-based stacking can misread occlusion; the
vision alternative needs comparison. CPU checks cover ordering, dynamic stacks, rejection,
schema recovery, safety routing and the unchanged-baseline transport adapter, and passed.
The combined runs still misinterpret throwing and pulling from under an object as ordinary
pick-and-place requests. Perception misses or mislabels the small white part and some saucer
fragments; stacked-object references and contradictory instructions remain unreliable.
The updated pipeline graph is exported in \texttt{figures/pipeline\_graph.pdf}.
\end{document}
